\begin{definition}[{\cite{MR2115083}}]\label{definitionalgebraicpartialaction}\index{Partial action!algebraic}
If $R$ is a commutative unital ring, $S$ is an inverse semigroup and $A$ is an associative algebra over $R$, a (non-degenerate) \emph{algebraic partial action} of $S$ on $A$ is a partial action $\alpha:S\to\mathcal{I}(A)$ such that:
\begin{enumerate}
\item[(i)] For all $s\in S$, $A_s$ is an ideal of $A$ and $\alpha_s:A_{s^*}\to A_s$ is an $R$-isomorphism;
\item[(ii)] $A=\operatorname{span}\bigcup_{s\in S}A_s$.
\end{enumerate}
Partial actions on rings are defined in the same way (e.g. by regarding rings as $\mathbb{Z}$-algebras in the natural way). All partial actions on algebras will be considered algebraic, in the sense described here.
\end{definition}

Similarly to the case of topological partial actions, Property (ii) can be substituted with
\begin{enumerate}
\item[(ii)'] $A=\operatorname{span}\bigcup_{e\in E(S)}A_e$;
\end{enumerate}
and this \emph{non-degeneracy} condition will always be assumed to hold, since we can always substitute $A$ with $\operatorname{span}\bigcup_{e\in E(S)}A_e$ if necessary.

Crossed products for partial actions of inverse semigroups are defined in the same way as crossed products of global actions (see \cite{MR2559043}), as long as we take appropriate care when dealing with \emph{partial} morphisms.

\begin{definition}\label{definitionLalpha}\nomenclature[P]{$\mathscr{L}(\alpha)$}{Definition \ref{definitionLalpha}}
Let $\alpha$ be a (algebraic) partial action of $S$ on a $R$-algebra $A$. We denote by $\mathscr{L}(\alpha)$ the $R$-module of all finite formal sums
\[\sum_{s\in S} a_s \delta_s,\qquad\text{where}\quad a_s\in A_s\quad\text{and}\quad\delta_s\quad\text{is a formal symbol}.\]

More precisely, $\mathscr{L}(\alpha)$ is the free $R$-module generated by formal elements $a_s\delta_s$, where $s\in S$ and $a_s\in A_s$, satisfying the relations
\[(a_s\delta_s)+(b_s\delta_s)=(a_s+b_s)\delta_s\qquad\text{and}\qquad \lambda(a_s\delta_s)=(\lambda a_s)\delta_s\]
for all $s\in S$, $a_s,b_s\in A_s$ and $\lambda\in R$.
\end{definition}

\begin{remark}
Alternatively, $\mathscr{L}(\alpha)$ coincides with the $R$-module of all finitely supported functions $f:S\to A$ which satisfy $f(s)\in A_s$ for all $s\in S$. One should be a somewhat careful with this interpretation, since we could think of $\delta_s$ as the characteristic function of $\left\{s\right\}$, which will not be an element of $\mathscr{L}(\alpha)$ in general. The index $s$ appearing in a term $a_s\delta_s$ will be useful to keep track of which ideal the element $a_s$ belongs to.
\end{remark}

We define a product on the generators of $\mathscr{L}(\alpha)$ by
\[(a_s \delta_s)(b_t \delta_t) = \alpha_{s}(\alpha_{s^*}(a_s) b_t) \delta_{st}.\]
%To check that this is well-defined, we use $\alpha_{s^*}(a_s)\in A_s$, which is an ideal, so $\alpha_{s^*}(a_s)b_t\in A_s$ as well and thus the coefficient $\alpha_{s}(\alpha_{s^*}(a_s) b_t)$ is defined. Since
%\[\alpha_s^*(\alpha_{s}(\alpha_{s^*}(a_s) b_t))=\alpha_{s^*}(a_s)b_t\in A_{s^*}\cap A_t\]
%then
%\[\alpha_{s}(\alpha_{s^*}(a_s) b_t)\in\alpha_{s^*}^{-1}(A_{s^*}\cap A_t)\subseteq A_{(t^*s^*)^*}=A_{st}\]
%and therefore the product is well-defined on the generators and can be extended to a bilinear map on $\mathscr{L}(\alpha)$.

\begin{example}
Let $\alpha$ be a global action of a group $G$ on an $R$-algebra $A$. Then, as a module, $\mathscr{L}(\alpha)$ coincides with the group algebra $A\rtimes_{\alpha} G$. Since all functions $\alpha_s$ are automorphisms of $A$, we can calculate the product on $\mathscr{L}(\alpha)$ as
\[(a_s\delta_s)(b_t\delta_t)=\alpha_s(\alpha_{s^{-1}}(a_s)b_t)\delta_st=\alpha_s(\alpha_{s^{-1}}(a_s))\alpha_s(b_t)\delta_{st}=a_s\alpha_s(b_t)\delta_{st}\]
so $\mathscr{L}(\alpha)$ is in fact the same as the group algebra $A\rtimes_{\alpha} G$.
\end{example}

In general, this product might make $\mathscr{L}(\alpha)$ a \emph{non-associative algebra} (\cite[Example~3.5]{MR2115083}). We will prove associativity in a particular case, which will always be satisfied for the algebras we will consider. For more genera conditions which guarantee associativity of $\mathscr{L}(\alpha)$, see 3.1 and 3.2 of \cite{MR2115083}.


\begin{definition}
We say that an algebra $A$
\begin{itemize}
\item is \emph{idempotent} if for all $a\in A$, there are finitely many elements $b_1,c_1,\ldots,b_n,c_n$ such that $a=\sum_i b_ic_i$;
\item has \emph{local units} if for every finite subset $F\subseteq A$, there exists $e\in A$ such that $e^2=e$ and $ef=fe=f$ for all $f\in F$, and call such $e$ a \emph{local unit} for $F$
\end{itemize}
\end{definition}

\begin{proposition}
Suppose that $\alpha=\left(\alpha_s:A_{s^*}\to A_s\right)_{s\in S}$ is an algebraic partial action of an inverse semigroup $S$ on an $R$-algebra $A$, and that every ideal $A_s$ is idempotent. Then $\mathscr{L}(\alpha)$ is associative.

In particular, if every ideal $A_s$ has local units then $\mathscr{L}(\alpha)$ is associative.
\end{proposition}
\begin{proof}
Of course, it is enough to prove that
\[(a_s\delta_s)\left((b_t\delta_t)(c_z\delta_z)\right)=\left((a_s\delta_s)(b_t\delta_t)\right)(c_z\delta_z)\]
for all $s,t,z\in S$ and $a_s\in A_s$, $b_t\in A_t$ and $c_z\in A_z$. Suppose initially that $b_t=b_t^1b_t^2$, where $b_t^i\in A_t$. Then
\begin{align*}
(a_s\delta_s)\left((b_t\delta_t)(c_z\delta_z)\right)\hspace{-3cm}&\hspace{3cm}=\alpha_s\left(\alpha_{s^*}(a_s)\alpha_t(\alpha_{t^*}(b_t)c_z)\right)\delta_{stz}=\alpha_s\left(\alpha_{s^*}(a_s)b_t^1\alpha_t(\alpha_{t^*}(b_t^2)c_z)\right)\delta_{stz}
\end{align*}
Since $\alpha_{s^*}b_t^1\in A_t$, we get
\begin{align*}
(a_s\delta_s)\left((b_t\delta_t)(c_z\delta_z)\right)\hspace{-0cm}&\hspace{0cm}=\alpha_s\left(\alpha_t(\alpha_{t^*}(\alpha_{s^*}(a_s)b_t^1)\alpha_t(\alpha_{t^*}(b_t^2)c_z))\right)\delta_{stz}\\
&=\alpha_s\left(\alpha_t\left[\alpha_{t^*}(\alpha_{s^*}(a_s)b_t^1\alpha_{t^*}(b_t^2)c_z\right]\right)\delta_{stz}\\
&=\alpha_{st}\left[\alpha_{t^*}(\alpha_{s^*}(a_s)b_t^1)\alpha_{t^*}(b_t^2)c_z\right]\delta_{stz}\\
&=\alpha_{st}\left[\alpha_{t^*}(\alpha_{s^*}(a_s)b_t^1b_t^2)c_z\right]\delta_{stz}\\
&=\alpha_{st}\left[\alpha_{t^*}(\alpha_{s^*}(a_s)b_t)c_z\right]\delta_{stz}
\end{align*}
Since $\alpha_{s^*}(a_s)b_t\in A_t$, we can rewrite this as
\begin{align*}
(a_s\delta_s)\left((b_t\delta_t)(c_z\delta_z)\right)&=\alpha_{st}\left[\alpha_{t^*}(\alpha_{s^*}(\alpha_s(\alpha_{s^*}(a_s)b_t)))c_z\right]\delta_{stz}\\
&=\alpha_{st}\left[\alpha_{(st)^*}(\alpha_s(\alpha_{s^*}(a_s)b_t))c_z\right]\delta_{stz}
\end{align*}
and this coincides precisely with the definition of $\left((a_s\delta_s)(b_t\delta_t)\right)(c_z\delta_z)$. Since every element of $A_t$ is a sum of products of elements of $A_t$ then we are done.\qedhere
\end{proof}

\begin{definition}\label{definitionNalpha}\nomenclature[P]{$\mathscr{N}(\alpha)$}{Definition \ref{definitionNalpha}}
Given a partial action $\alpha$ of an inverse semigroup $S$ on an algebra $A$, we let $\mathscr{N}(\alpha)$ be the additive subgroup of $\mathscr{L}(\alpha)$ generated by all elements of the form
\[a_r\delta_r - a_r\delta_s,\qquad\text{where}\quad r\leq s\quad\text{and}\quad a_r\in A_r.\] 
\end{definition}

\begin{proposition}
If $\alpha$ is an action of $S$ on an algebra $A$, then $\mathscr{N}(\alpha)$ is a two-sided ideal of $\mathscr{L}(\alpha)$.
\end{proposition}
\begin{proof}
If $a_r\delta_r-a_r\delta_s$, where $a\in A_r$ and $r\leq s$ and $\lambda\in R$, then
\[\lambda(a_r\delta_r-a_r\delta_s)=(\lambda a_r)\delta_r-(\lambda a_r)\delta_s\]
so $\mathscr{N}(\alpha)$ is a submodule of $\mathscr{L}(\alpha)$. If $b_z\delta_z\in\mathscr{L}$, then
\[(b_z\delta_z)(a_r\delta_r-a_r\delta_s)=\alpha_z(\alpha_{z^*}(b_z)a_r)\delta_{zr}-\alpha_z(\alpha_{z^*}(b_z)a_r)\delta_{zs}\]
and since $zr\leq zs$ then $\mathscr{N}(\alpha)$ is a left ideal of $\mathscr{L}(\alpha)$, and similarly it is also a right ideal.\qedhere
\end{proof}

\begin{definition}\label{definitioncrossedproduct}\nomenclature[P]{$A\rtimes_\alpha S$}{Crossed product}
Given an (algebraic) partial action $\alpha$ of $S$ on an algebra $A$, the (algebraic) \emph{crossed product} $A\rtimes_\alpha S$ is the quotient algebra
\[A\rtimes_\alpha S=\mathscr{L}(\alpha)/\mathscr{N}(\alpha)\]
and as usual, we will simply write $\mathscr{L}$, $\mathscr{N}$ and $A\rtimes S$ when there is no risk of confusion.% Partial crossed product rings are defined in the same manner, by regarding rings as $\mathbb{Z}$-algebras.

The class of an element $x\in\mathscr{L}(\alpha)$ in $\Skew$ will be denoted by $\overline{x}$.
\end{definition}

\begin{remarkn}\label{remarkpresentationofalgebraiccrossedproduct}
The difference from $\mathscr{L}$ to $A\rtimes S$ is that $A\rtimes S$ respects the order structure of $S$: Using the presentation of $\mathscr{L}(\alpha)$ given in Definition \ref{definitionLalpha}, we obtain the following presentation of $A\rtimes S$: It is the free module generated by symbols $\overline{a_s\delta_s}$ where $s\in S$ and $a_s\in A_s$, and satisfying the following relations:
\begin{itemize}
\item $\overline{a_s\delta_s}+\overline{b_s\delta_s}=\overline{(a_s+ b_s)\delta_s}$, whenever $a_s,b_s\in A_s$;
\item $\lambda\left(\overline{a_s\delta_s}\right)=\overline{(\lambda a_s)\delta_s}$, whenever $a_s\in A_s$ and $\lambda\in R$;
\item $\overline{a_r\delta_r}=\overline{a_r\delta_s}$, whenever $a_r\in A_r$ and $r\leq s$.
\end{itemize}
\end{remarkn}

These are sometimes called ``\emph{partial} crossed products'' or ``\emph{partial skew inverse semigroup algebras}'' (e.g.\ in \cite{MR3669934} and \cite{MR3743184}, respectively), while the name ``crossed product'' (or ``skew inverse semigroup algebra'') are reserved for global actions. Instead, we adopt the simplest nomenclature.

Note that, since $\mathscr{N}$ is defined, in principle, simply as an additive subgroup of $\mathscr{L}$, then every map $\phi:\mathscr{L}\to B$, where $B$ is any algebra (or even simply a group), which satisfies $\phi(a_r\delta_r)=\phi(a_r\delta_s)$ whenever $s\leq r$ factors through $A\rtimes S$.

\begin{definition}\label{definitiondiagonalcrossedproduct}
The \emph{diagonal subalgebra} of a partial crossed product $A\rtimes S$ is the subalgebra generated by elements of the form $\overline{a\delta_e}$, where $e\in E(S)$ and $a\in A_e$.
\end{definition}

\begin{proposition}
Suppose $\alpha$ is a (non-degenerate), algebraic partial action of $S$ on $A$, and assume that
\begin{enumerate}
\item[(i)] $A$ has local units; and
\item[(ii)] every ideal $A_e$, $e\in E(S)$, has local units.
\end{enumerate}
Then the map $\iota:A\to A\rtimes S$, given by $a\mapsto \sum_{i=1}^n \overline{a_i\delta_{e_i}}$, where $e_i\in E(S)$ and $a_i\in A_{e_i}$ are chosen such that $a=\sum_{i=1}^n a_i$, is a well-defined injective algebra morphism.
\end{proposition}
\begin{proof}
We need to prove that the map $\iota$ is well-defined, that is, if $e_1,\ldots,e_n$ are idempotents and $a_i\in A_{e_i}$ are such that $\sum_{i=1}^n a_i=0$ in $A$, then $\sum_{i=1}^n\overline{a_i\delta_{e_i}}=0$ in $A\rtimes S$. We proceed by induction. If $n=1$ this is trivial, so suppose that the statement is valid for sums up to $n$ elements.

Suppose that $a+\sum_{i=1}^n a_i=0$, where $a\in A_e$ and $a_i\in A_{e_i}$ for certain $e,e_i\in E(S)$. Let $u\in A$ be a local unit for $\left\{a,a_1,\ldots,a_n\right\}$, and $v\in A_e$ be a local unit for $\left\{a\right\}$. Then
\[0=(u-v)a=-\sum_{i=1}^n(u-v)a_i,\]
and by the induction hypothesis,
\[0=\sum_{i=1}^n\overline{\left((u-v)a_i\right)\delta_{e_i}}=\sum_{i=1}^n\overline{a_i\delta_{e_i}}-\sum_{i=1}^n\overline{(va_i)\delta_{e_i}}\ntag\label{equationdiagonalsubalgebraforcrossedproducts1}.\]
However, $va_i\in A_{e_i}\cap A_e\subseteq A_{ee_i}$, so
\[\sum_{i=1}^n\overline{(va_i)\delta_{e_i}}=\sum_{i=1}^n\overline{(va_i)\delta_{ee_i}}=\sum_{i=1}^n\overline{(va_i)\delta_e}=\overline{\left(v\sum_{i=1}^n va_i\right)\delta_e}=\overline{-va\delta_e}=\overline{-a\delta_e},\]
and the previous equation implies that $\overline{a\delta_e+\sum_{i=1}^na_i\delta_{e_i}}=0$, as we wanted.

It is readily seen that $\iota$ is an module morphism, and to check that it is an algebra morphism, take generating elements of $A$ of the form $a\in A_{e}$, $b\in A_{f}$, where $e,f\in E(S)$. Then $ab\in A_{ef}$, and
\[\iota(a)\iota(b)=\overline{(a\delta_e)(b\delta_f)}=\overline{\alpha_e(\alpha_{e^*}(a)b)\delta_{ef}}=\overline{(ab)\delta_{ef}}=\iota(ab).\]

Now consider the module morphism $\tau:A\rtimes S\to A$ satisfying $\tau(\overline{a_s\delta_s})=a_s$, which exists because of the presentation of $A\rtimes S$ given in Remark \ref{remarkpresentationofalgebraiccrossedproduct}. Then as $\tau\circ\iota$ is the identity of $A$, $\iota$ is injective.\qedhere
\end{proof}

\begin{remark}
Note that the map $\tau$ above is not an algebra morphism in general. For example, let $S=\langle g|g^2=1\rangle$ is the cyclic group of order $2$ and $A=\mathbb{C}^2$ with the pointwise structure of $\mathbb{C}$-algebra. Define an action $\alpha$ of $S$ on $A$ as
\[\alpha_g(a,b)=(b,a),\qquad\alpha_1(a,b)=(a,b)\]
for all $(a,b)\in A$. Let $x=(1,0)\delta_g$ and $y=(0,1)\delta_1$. Then $xy=x$, so
\[\tau(xy)=\tau(x)=(1,0)\]
however,
\[\tau(x)\tau(y)=(1,0)(0,1)=(0,0).\]
\end{remark}

We will now describe a universal property for crossed products. Intuitively, $A\rtimes_\alpha S$ is the algebra generated by $A$ and $S$ in a way that the partial action $\alpha$ becomes conjugation by elements of $S$. However in order for $S$ to be included in $A\rtimes_\alpha S$, we need that all ideals $A_s$ have units.

\begin{definition}\label{definitioncovariantrepresentation}\index{Covariant representation}
Let $\alpha$ be a partial action of $S$ on an algebra $A$. A \emph{covariant representation} of $\alpha$ consists of an algebra $B$ and a pair $(\theta,\phi)$, where $\theta:S\to B$ is a partial morphism of $S$ to the multiplicative semigroup of $B$, and $\phi:A\to B$ is an algebra morphism, satisfying the \emph{covariance} condition
\[\phi(\alpha_s(a_{s^*}))=\theta(s)\phi(a_{s^*})\theta(s^*)\qquad\text{for all }s\in S\text{ and }a_{s^*}\in A_{s^*}.\]
We say that $(\theta,\phi)$ is \emph{non-degenerate} if $\theta(s)\theta(s^*)\in\phi(A_s)$ for all $s\in S$ (and \emph{degenerate} otherwise).
\end{definition}

\begin{example}
Let $X$ be a non-compact, locally compact Hausdorff space, and $S$ be the semigroup under intersection of clopen subsets of $X$. Let $A=C_c(X)$ be the (real or complex) algebra of compactly supported continuous functions on $X$. For each $U\in S$, let $A_U=\left\{f\in C_c(X):\supp(f)\subseteq U\right\}$, and $\alpha_U$ be the identity on $A_U$. Then $\alpha$ is a global action on $A$.

Let $B=C_b(X)$ be the algebra of continuous bounded functions on $X$. Define a covariant representation $(\theta,\phi)$ of $\alpha$ on $B$ by letting $\phi:A\hookrightarrow B$ be the inclusion, and $\theta(U)$ be the characteristic function of $U\in S$. Then $(\theta,\phi)$ is degenerate, because $\theta(X)\theta(X)=\theta(X)$ is the constant function at $1$, which does not belong to $\phi(A)$.
\end{example}

\begin{lemma}\label{lemmaforuniversalpropertyofcrossedproducts}
Let $(\theta,\phi)$ be a covariant representation of an algebraic partial action $\alpha:S\to\mathcal{I}(A)$.
\begin{enumerate}
\item[(a)] If $a_s\in A_s$ then
\[\phi(a_s)=\theta(s)\theta(s^*)\phi(a_s)=\phi(a_s)\theta(s)\theta(s^*).\ntag\label{equationlemmaforuniversalpropertyofcrossedproductsitema}\]
\item[(b)] If $s\leq t$ and $a_s\in A_s$ then $\phi(a_s)\theta(s)=\phi(a_s)\theta(t)$.
\item[(c)] If $A_s$ has a unit $1_s$ and $(\theta,\phi)$ is non-degenerate, then
\[\phi(1_s)\theta(s)=\theta(s)\qquad\text{and}\qquad\theta(s^*)\phi(1_s)=\theta(s^*).\]
\end{enumerate}
\end{lemma}
\begin{proof}
\begin{enumerate}
\item[(a)] As $a_s=\alpha_s(\alpha_{s^*}(a_s))$, by covariance we get
\[\phi(a_s)=\theta(s)\theta(s^*)\phi(a_s)\theta(s)\theta(s^*).\]
The partial morphism properties of $\theta$ imply
\[\theta(s)\theta(s^*)\theta(s)\theta(s^*)=\theta(s)\theta(s^*),\]
and Equation \ref{equationlemmaforuniversalpropertyofcrossedproductsitema} follows from these two equalities.
\item[(b)] Item (a) and $\theta$ being a partial morphism imply that, if $s\leq t$,
\begin{align*}
\phi(a_s)\theta(t)&=\phi(a_s)\theta(s)\theta(s^*)\theta(t)=\phi(a_s)\theta(s)\theta(s^*t)\\
&=\phi(a_s)\theta(s)\theta(s^*s)=\phi(a_s)\theta(s)\theta(s^*)\theta(s)=\phi(a_s)\theta(s).
\end{align*}
\item[(c)] Since $(\theta,\phi)$ is non-degenerate, then $\theta(s)\theta(s^*)\in\phi(A_s)$, so
\[\theta(s)=\theta(s)\theta(s^*)\theta(s)=\phi(1_s)\theta(s)\theta(s^*)\theta(s)\overset{(a)}{=} \phi(1_s)\theta(s)\]
and the other equality is similar.\qedhere
\end{enumerate}
\end{proof}

\begin{theorem}\label{theoremuniversalpropertyofcrossedproducts}
Suppose that $\alpha$ is an algebraic partial action of an inverse semigroup $S$ on an algebra $A$ with local units, and such that $A_s$ has a unit $1_s$ for all $s\in S$. Define
\[\sigma:S\to A\rtimes S\quad\text{by}\quad \sigma(s)=\overline{1_s\delta_s}.\]
Let $\iota:A\to A\rtimes S$ be the embedding of $A$ as the diagonal subalgebra of $A\rtimes S$. Then $(\sigma,\iota)$ is a universal non-degenerate covariant representations of $\alpha$ in the following sense:
\begin{enumerate}
\item[(i)] $(\sigma,\iota)$ is a non-degenerate covariant representation;
\item[(ii)] If $(\theta,\phi)$ is any other non-degenerate covariant representation of $\alpha$ on an algebra $B$, then there exists a unique algebra morphism $\Phi:A\rtimes S\to B$ such that
\[\phi=\Phi\circ\iota\quad\text{and}\quad\theta=\Phi\circ\sigma.\]
\end{enumerate}
\end{theorem}
\begin{proof}
First we prove that $\sigma$ is a partial morphism. Given $s,t\in S$, we have $\alpha_s(1_{s^*})=1_s$ because $\alpha_s:A_{s^*}\to A_s$ is an algebra isomorphism, and it follows that
%\[\sigma(s)\overline{(a_t\delta_t)}=\overline{(1_s\delta_s)(a_t\delta_t)}=\overline{\alpha_s(\alpha_{s^*}(1_s)a_t)\delta_{st}}=\overline{\alpha_s(1_{s^*}a_t)\delta_{st}}\]
%and in particular
\[\sigma(s)\sigma(s^*)=\overline{(1_s)\delta_s(1_{s^*})\delta_{s^*}}=\overline{1_s\delta_{ss^*}}.\ntag\label{equationnondegeneracyuniversalpropertycrossedproduct}\]
\begin{enumerate}
\item[(PM1)]  Using Equation \ref{equationnondegeneracyuniversalpropertycrossedproduct}, we compute
\[\sigma(s)\sigma(s^*)\sigma(s)=\overline{1_s\delta_s}\overline{1_{s^*}\delta_{s^*s}}=\overline{1_s\delta_{ss^*s}}=\sigma(s).\]
\item[(PM2)] We have
\[\sigma(s)\sigma(t)\sigma(t^*)=\overline{\alpha_s(1_{s^*}1_t)\delta_{stt^*}}\quad\text{and}\quad\sigma(st)\sigma(t^*)=\overline{\alpha_{st}(1_{(st)^*}1_{t^*})\delta_{stt^*}}.\]

On one hand, $1_{s^*}1_t$ is the unit of $A_{s^*}\cap A_t$, so $\alpha_s(1_{s^*}1_t)$ is the unit of $\alpha_s(A_{s^*}\cap A_t)$.

On the other hand, $1_{(st)^*}1_{t^*}$ is the unit of $A_{(st)^*}\cap A_{t^*}$, so $\alpha_{st}(1_{(st)^*}1_{t^*})$ is the unit of $\alpha_{st}(A_{(st)^*}\cap A_{t^*})$. Proposition \ref{propositioncomparingdomainspartialaction}(b) and (c) imply
\begin{align*}
\alpha_s(A_{s^*}\cap A_t)&\overset{\ref{propositioncomparingdomainspartialaction}(b)}{=}A_s\cap A_{st}\overset{\ref{propositioncomparingdomainspartialaction}(c)}{=}A_{(st)(st)^*s}\cap A_{st}=A_{stt^*}\cap A_{st}\\\
&\overset{\ref{propositioncomparingdomainspartialaction}(b)}{=}\alpha_{st}(A_{(st)^*}\cap A_{t^*}),
\end{align*}
and we conclude that $\sigma(s)\sigma(t)\sigma(t^*)=\sigma(st)\sigma(t^*)$.
\item[(PM3)] As in the previous item, we compute
\[\sigma(s^*)\sigma(s)\sigma(t)=\overline{1_{s^*}1_t\delta_{s^*st}}\quad\text{and}\quad\sigma(s^*)\sigma(st)=\overline{\alpha_{s^*}(1_s1_{st})\delta_{s^*st}}\]
and note that both coefficients are the units of
\[\alpha_{s^*}(A_s\cap A_{st})=A_{s^*}\cap A_{s^*st}=A_{s^*}\cap A_t\]
again using Proposition \ref{propositioncomparingdomainspartialaction}(b) and (c).
\end{enumerate}

We already know that $\iota$ is an algebra morphism. Using the definition of the product in $A\rtimes S$ and the fact that $\alpha_s(1_{s^*})=1_s$ we immediately obtain covariance of $(\sigma,\iota)$, and non-degeneracy follows from Equation (\ref{equationnondegeneracyuniversalpropertycrossedproduct}).% If $s\in S$ and $a_s\in A_{s^*}\subseteq A_{s^*s}$, then
%\begin{align*}
%\sigma(s)\iota(a_{s^*})\sigma(s^*)&=\overline{(1_s\delta_s)(a_{s^*}\delta_{s^*s})(1_{s^*}\delta_{s^*})}=\overline{(\alpha_{s}(1_{s^*}a_{s^*})\delta_{ss^*s})(1_{s^*}\delta_{s^*})}\\
%&=\overline{(\alpha_s(a_{s^*})\delta_s)(1_{s^*}\delta_{s^*})}=\overline{\alpha_s(\alpha_{s^*}(\alpha_s(a_{s^*}))1_{s^*})\delta_{ss^*}}\\
%&=\overline{\alpha_s(a_{s^*})\delta_{ss^*}}
%\end{align*}
%and since $\alpha_s(a_{s^*})\in A_s\subseteq A_{ss^*}$, we conclude that
%\[\sigma(s)\iota(a_{s^*})\sigma(s^*)=\overline{\iota(\alpha_s(a_{s^*}))}.\]
%
%To check that $(\sigma,\iota)$ is non-degenerate, simply notice that
%\[\sigma(s)\sigma(s^*)=\overline{1_s\delta_{ss^*}}=\iota(1_s)\in\iota(A_s).\]

As for the universal property, suppose $(\theta,\phi)$ is a covariant representation of $\alpha$ on an algebra $B$. Using the presentation of $A\rtimes S$ in Remark \ref{remarkpresentationofalgebraiccrossedproduct} and Lemma \ref{lemmaforuniversalpropertyofcrossedproducts}(b), we can define a module morphism $\Phi:A\rtimes S\to B$ satisfying $\Phi(\overline{a_s\delta_s})=\phi(a_s)\theta(s)$ for all $a_s\in A_s$ ($s\in S$).

Let us check that $\Phi$ is an algebra morphism: Let $\overline{a_s\delta_s},\overline{b_t\delta_t}$ be generators of $A\rtimes S$. Then
\[\Phi(\overline{(a_s\delta_s)(b_t\delta_t)})=\Phi(\overline{\alpha_s(\alpha_{s^*}(a_s)b_t)\delta_{st}})=\phi(\overline{\alpha_s(\alpha_{s^*}(a_s)b_t)}\theta(st))\]
and using covariance and $\phi$ being multiplicative,
\[\Phi(\overline{(a_s\delta_s)(b_t\delta_t)})=\theta(s)\theta(s^*)\phi(a_s)\theta(s)\phi(b_t)\theta(s^*)\theta(st)\ntag\label{equationtheoremuniversalpropertyofcrossedproducts1}\]
On the other hand
\[\Phi(\overline{a_s\delta_s})\Phi(\overline{b_t\delta_t})=\phi(a_s)\theta(s)\phi(b_t)\theta(t)\ntag\label{equationtheoremuniversalpropertyofcrossedproducts1.5}\]
so we will modify the right-hand term of Equation (\ref{equationtheoremuniversalpropertyofcrossedproducts1}) to obtain the right-hand term of (\ref{equationtheoremuniversalpropertyofcrossedproducts1.5}).

From Lemma \ref{lemmaforuniversalpropertyofcrossedproducts}(a)
\[\theta(s)\theta(s^*)\phi(a_s)=\phi(a_s).\ntag\label{equationtheoremuniversalpropertyofcrossedproducts2}\]

By Lemma \ref{lemmaforuniversalpropertyofcrossedproducts}(c) and (a), (and the inclusion $1_{s^*}b_t\in A_{s^*}$), we obtain
\begin{align*}
\theta(s)\phi(b_t)\theta(s^*)\theta(st)&=\theta(s)\phi(1_{s^*}b_t)\theta(s^*)\theta(s)\theta(t)=\theta(s)\phi(1_{s^*}b_t)\theta(t)\\
&=\theta(s)\phi(b_t)\theta(t)\ntag\label{equationtheoremuniversalpropertyofcrossedproducts3}
\end{align*}
Using Equations (\ref{equationtheoremuniversalpropertyofcrossedproducts2}), (\ref{equationtheoremuniversalpropertyofcrossedproducts3}), we conclude from (\ref{equationtheoremuniversalpropertyofcrossedproducts1}) and (\ref{equationtheoremuniversalpropertyofcrossedproducts1.5}) that $\Phi$ is an algebra morphism.

Given $e\in E(S)$ and $a_e\in A_e$, Lemma \ref{lemmaforuniversalpropertyofcrossedproducts}(a) and $\theta(e)$ being idempotent imply that
\[\Phi(\iota(a_e))=\phi(a_e)\theta(e)=\phi(a_e)\theta(e)\theta(e^*)=\phi(a_e)\]
and so $\Phi\circ\iota=\phi$ on $A_e$ and thus on all of $A$. The fact that $\Phi\circ\sigma=\theta$ follows from Lemma \ref{lemmaforuniversalpropertyofcrossedproducts}(c).

Finally, uniqueness of $\Phi$ with the properties in (ii) is a consequence of $\overline{a_s\delta_s}$ being equal to $\iota(a_s)\sigma(\delta_s)$ for all $s\in S$ and $a_s\in A_s$.\qedhere
\end{proof}

%\begin{theorem}
%Suppose $\alpha$ is a partial action of $S$ on an algebra $A$ such that $A_s$ has has local units for all $s$. Let $B$ be an ideal of $A$, and $\beta$ the compression of $\alpha$ to $B$, that is $B_s=A_s\cap B$ and $\beta_s$ is the restriction of $\alpha_s$ to $B_{s^*}$.
%
%Then the map $B\rtimes_\beta S$ is isomorphic to an ideal of $A\rtimes_\alpha S$ (in the ``obvious way'').
%\end{theorem}
%\begin{proof}
%Since $\mathscr{L}(\beta)$ is an ideal of $\mathscr{L}(\alpha)$, we can consider the map $\Psi:\mathscr{L}(\beta)\to A\rtimes_\alpha S$, $\Psi(b_t\delta_t)=\overline{b_t\delta_}$. We are done if we prove that $\ker\Psi=\mathscr{N}(\beta)$, because in this case $\Psi$ factors though an injective morphism of $B\rtimes_\beta S$.
%
%Consider on $\mathscr{L}(\alpha)$ the $B$-module structure given by $b(a_t\delta_t)=(bat_)\delta_t$, ($b\in B, a_t\in A_t$). 
%
%So given $x\in\ker\Psi=\mathscr{N}(\alpha)\cap\mathscr{L}(\beta)$, we can write $x=\sum_{i=1}^n a_i\delta_{r_i}-a_i\delta_{s_i}$, where $r_i\leq s_i$, $a_i\in A_{r_i}$. But since $x\in\mathscr{L}(\beta)$ and $B$ has local units, we can find $u\in B$ with $ux=x$, so
%\[x=ux=\sum_{i=1}^n(ua_i)\delta_{r_i}-(ua_i)\delta_{s_i}\in\mathscr{N}(\beta)\]
%because $ua_i\in B\cap A_{r_i}=B_{r_i}$ for all $i$.
%\end{proof}